% ======================================================================
\section{Using the framework}\label{sec:framework}
% ======================================================================

This section builds a complete agent from nothing. The running example is a
solar-powered borehole pump serving a village: it has a tank-level sensor, a
turbidity probe, a battery monitor and a flow meter, plus a text field where a
caretaker types what they see. There is no reliable network. The full file is
\file{examples/pump\_station.py} --- run it with
\code{python examples/pump\_station.py}.

\begin{keybox}[The shape of every CONCLAVE program]
\textbf{tools} $\rightarrow$ \textbf{training observations} $\rightarrow$
\textbf{featuriser} $\rightarrow$ \textbf{brain} $\rightarrow$
\textbf{gate + constraints} $\rightarrow$ \textbf{agent}.\\[2mm]
Six steps, in that order, always. The rest is detail.
\end{keybox}

\subsection{Step 1 --- declare what the agent can do}

A tool is a Python function plus a schema. The \code{@tool} decorator infers the
schema from the type annotations, the way a typed route converter does:

\begin{codebox}[Tools]
\begin{lstlisting}[style=py]
from conclave.tools.base import tool
from conclave import ToolRegistry

@tool(description="Start the borehole pump to refill the tank")
def start_pump(minutes: int = 30) -> str:
    return f"pump running {minutes} min"

@tool(description="Open the chlorine doser to treat the tank", irreversible=True)
def dose_chlorine(ml: int = 50) -> str:
    return f"dosed {ml} ml"

@tool(description="Send an SMS to the maintenance technician")
def call_technician(reason: str = "unknown") -> str:
    return f"SMS sent: {reason}"

registry = ToolRegistry([start_pump, dose_chlorine, call_technician])
\end{lstlisting}
\end{codebox}

\noindent Three arguments to \code{@tool} carry real weight:

\begin{center}
\renewcommand{\arraystretch}{1.25}
\begin{tabular}{@{}>{\raggedright\arraybackslash}p{3.0cm}
                 >{\raggedright\arraybackslash}p{12.2cm}@{}}
\toprule
\rowcolor{conclavebg}
{\sffamily\bfseries\color{conclaveblue}Argument} &
{\sffamily\bfseries\color{conclaveblue}Effect} \\
\midrule
\code{description} & Becomes matchable text for \code{ToolRouter}. Write it the
way a user would describe the need, not the way you named the function. \\
\code{irreversible} & Marks the tool for the constraint layer and lets the
planner exclude it. Dosing chemicals is irreversible; sending an SMS is not. \\
\code{preconditions} / \code{effects} & Predicate lists the planner searches
over. Leave empty if you are not planning. \\
\bottomrule
\end{tabular}
\end{center}

\subsection{Step 2 --- observations}

An \code{Observation} carries numeric features, optional free text, and metadata.
Either modality may be empty; the featuriser handles whichever is present.

\begin{codebox}[Observations]
\begin{lstlisting}[style=py]
from conclave import Observation

Observation(
    features={"tank_level_pct": 11.0, "turbidity_ntu": 1.2,
              "battery_volts": 12.9, "flow_lpm": 0.0},
    text="tank nearly empty this morning",
)
\end{lstlisting}
\end{codebox}

Training data is a list of observations and a matching list of action names.
In the example these are generated synthetically, five scenarios $\times$ 120
samples; in a real deployment they are logged decisions from whoever runs the
station today.

\begin{warnbox}[Two constraints on your training data]
\begin{itemize}[leftmargin=*,itemsep=2pt,topsep=2pt]
  \item \code{EnsembleBrain} needs \textbf{at least two classes}, and enough
        examples per class for its internal cross-validation. A handful of rows
        will raise.
  \item Class names become action names. They must match tool names exactly, or
        \code{Agent} will decide on an action that is not in the registry and
        silently do nothing.
\end{itemize}
\end{warnbox}

\subsection{Step 3 --- featurise under a budget}

\begin{codebox}[Featurisation]
\begin{lstlisting}[style=py]
from conclave.brain.featurize import (BudgetedTextFeaturizer, FusedFeaturizer,
                                   TabularFeaturizer)

fz = FusedFeaturizer(
    text=BudgetedTextFeaturizer(n_word_features=2**12, n_char_features=2**12),
    tabular=TabularFeaturizer(),
).fit(observations)

X = fz.transform(observations)
\end{lstlisting}
\end{codebox}

\begin{keybox}[The feature budget is the deployment knob]
\code{BudgetedTextFeaturizer} uses the hashing trick, so it stores
\textbf{no vocabulary at all} --- its size is fixed when you construct it and
does not grow with your corpus. \code{n\_word\_features} and
\code{n\_char\_features} therefore set the model size directly:
$2^{12}$ gives a 2.6\,MB artefact, $2^{16}$ gives 36\,MB. Table~III of the
paper shows what that costs in accuracy. \S\ref{sec:math} explains why the
hashing trick works and what it costs statistically.
\end{keybox}

\subsection{Step 4 --- the brain}

\begin{codebox}[The committee]
\begin{lstlisting}[style=py]
from conclave import EnsembleBrain

brain = EnsembleBrain().fit(X, labels)

print(list(brain.members_))                        # ['logreg', 'sgd_mh', 'nb']
print(brain.weights_.round(3))                     # [0.333 0.333 0.333]
print(brain.footprint_bytes()["_total"] / 1000)    # 24.2  (KB)
\end{lstlisting}
\end{codebox}

\noindent Twenty-four kilobytes. That is the entire decision-making apparatus
for this agent. \S\ref{sec:models} covers what is inside and how to change it.

\subsection{Step 5 --- the gate and the safety layer}

\begin{codebox}[Uncertainty gate and hard constraints]
\begin{lstlisting}[style=py]
from conclave import UncertaintyGate
from conclave.planning.rules import ConstraintSet, require_confidence

gate = UncertaintyGate(epistemic_max=0.30,   # committee disagreement ceiling
                       aleatoric_max=0.75,   # genuine-ambiguity ceiling
                       confidence_min=0.45,
                       margin_min=0.05,
                       allow_gather=False)   # no cheap extra sensing here

# No statistical model may dose chemicals on a weak majority.
constraints = ConstraintSet([require_confidence(0.80, ["dose_chlorine"])])
\end{lstlisting}
\end{codebox}

\begin{keybox}[The constraint layer sits \emph{above} the learned policy]
A violated constraint does not adjust a probability --- it converts the action
into a recorded \code{ESCALATE}. This is deliberate: no statistical model should
be able to authorise an irreversible act on its own, however confident it
claims to be.
\end{keybox}

\noindent Rather than hand-tuning those four thresholds, you can calibrate them
against a target precision on held-out data:

\begin{codebox}[Calibrating instead of guessing]
\begin{lstlisting}[style=py]
gate = brain.calibrate_gate(X_val, y_val, target_precision=0.90)
print(gate.calibration_)
# {'target_precision': 0.9,
#  'achieved': {'coverage': 0.52, 'precision': 0.906}, ...}
\end{lstlisting}
\end{codebox}

\noindent Read that as: \emph{acting on 52\% of decisions and deferring the rest
buys 90.6\% accuracy on the ones acted upon}. \S\ref{sec:selective} explains the
selective-prediction machinery behind it.

\subsection{Step 6 --- assemble and decide}

\begin{codebox}[The agent]
\begin{lstlisting}[style=py]
from conclave import Agent, EpisodicMemory

agent = Agent(brain=brain, featurizer=fz, registry=registry, gate=gate,
              constraints=constraints, memory=EpisodicMemory(capacity=500))

d = agent.decide(Observation(features=..., text="very cloudy water"))
print(d.action, d.confidence, d.disagreement, d.gate)
\end{lstlisting}
\end{codebox}

\noindent Running the four probe observations from the example file gives:

\begin{codebox}[Actual output of examples/pump\_station.py]
\begin{lstlisting}[style=out]
observation                            action             conf  disag      gate
tank nearly empty this morning         <escalate>         0.65   0.40  escalate
very cloudy water since the storm      dose_chlorine()    1.00   0.00       act
no flow and the battery looks flat     call_technician()  0.78   0.19       act
strange smell from the wellhead ...    <escalate>         0.46   0.35  escalate
\end{lstlisting}
\end{codebox}

All four rows are worth reading carefully, because two of them are the agent
working and two are the agent \emph{knowing it should not act}:

\begin{center}
\renewcommand{\arraystretch}{1.3}
\begin{tabular}{@{}>{\raggedright\arraybackslash}p{4.0cm}
                 >{\raggedright\arraybackslash}p{11.2cm}@{}}
\toprule
\rowcolor{conclavebg}
{\sffamily\bfseries\color{conclaveblue}Row} &
{\sffamily\bfseries\color{conclaveblue}What actually happened} \\
\midrule
cloudy water & Textbook case. Every member agrees (disagreement 0.00), confidence
1.00 clears the 0.80 chlorine constraint, the tool fires. \\
no flow, flat battery & Correct call at 0.78 confidence with mild disagreement.
Note it would \emph{not} have been allowed to dose chlorine at this confidence. \\
strange smell & Nothing like it in training. The committee splits (0.35), the gate
escalates rather than guessing. \textbf{This is the behaviour the whole design
exists to produce.} \\
tank nearly empty & An honest failure. Training said ``tank low overnight draw
down''; the probe says ``nearly empty''. Hashed $n$-grams have no idea those mean
the same thing, so the members split and it escalates. \\
\bottomrule
\end{tabular}
\end{center}

\begin{warnbox}[That last row is a real limitation, not a bug]
CONCLAVE's text features are lexical. They have no compositional semantics, so
paraphrase, negation and synonymy hurt it more than they hurt a transformer.
The paper states this in its limitations; the example reproduces it rather than
hiding it. If your text varies a lot in phrasing, either widen your training
vocabulary or lean harder on the numeric features, which have no such problem.
\end{warnbox}

\subsection{What a \code{Decision} carries}

Every decision is fully auditable --- this matters when an agent is making
consequential calls unsupervised:

\begin{center}
\renewcommand{\arraystretch}{1.22}
\begin{tabular}{@{}>{\raggedright\arraybackslash}p{3.6cm}
                 >{\raggedright\arraybackslash}p{11.6cm}@{}}
\toprule
\rowcolor{conclavebg}
{\sffamily\bfseries\color{conclaveblue}Field} &
{\sffamily\bfseries\color{conclaveblue}Meaning} \\
\midrule
\code{action}        & What it decided: a tool call, or abstain / gather / escalate \\
\code{posterior}     & Probability over every action \\
\code{member\_votes} & What \emph{each} committee member believed, separately \\
\code{confidence}    & Largest posterior probability \\
\code{disagreement}  & Normalised epistemic uncertainty --- the spread of opinion \\
\code{entropy}       & Normalised total uncertainty \\
\code{margin}        & Top-1 minus top-2 probability \\
\code{gate}          & Which branch fired, and therefore why \\
\code{latency\_ms}   & Wall-clock cost of this decision \\
\bottomrule
\end{tabular}
\end{center}

\subsection{Running a sequential environment}

\code{agent.decide()} handles one observation. For a control loop, \code{agent.run()}
drives any environment exposing \code{reset()} and \code{step(action)}:

\begin{codebox}[Episodes]
\begin{lstlisting}[style=py]
from conclave.envs.disaster import DisasterConfig, FloodResponseEnv

episode = agent.run(FloodResponseEnv(DisasterConfig(outage_prob=0.3)),
                    max_steps=40)

print(episode.total_reward, episode.n_escalations)
print(episode.to_dict())          # JSON-serialisable audit trail
\end{lstlisting}
\end{codebox}

\noindent With \code{memory} attached and \code{learn\_online=True} (the default),
successful steps are written back into episodic memory as cases, so the agent
adapts within a deployment without retraining.

\subsection{Tool routing for open vocabularies}

\code{EnsembleBrain} classifies into a \emph{fixed} action set. When tools come
and go at runtime, use \code{ToolRouter} instead: it scores \emph{(request, tool)}
pairs, so a tool registered after training is scored by the same features as
any other.

\begin{codebox}[Routing]
\begin{lstlisting}[style=py]
from conclave.brain.router import ToolRouter

router = ToolRouter(abstain=True, abstain_threshold=0.35).fit(
    queries, candidate_sets, targets)

result = router.route("people stuck on a roof in sector 4, send help",
                      registry.specs())
print(result.tool, result.scores, result.latency_ms)
\end{lstlisting}
\end{codebox}

\noindent \code{result.tool} is \code{None} when nothing scores above the
abstention threshold --- which is the correct answer more often than agent
frameworks like to admit. Try \code{python demo.py route} and give it something
outside its competence; it will show you flat scores rather than a confident
mistake.
